\hypertarget{appendix-215ac7099a82}{}\label{referencestart:appendix-215ac7099a82}
\section{Minkowski sums and support mappings}\label{appendix-215ac7099a82-minkowski-sums-and-support-mappings}

\textbf{Curriculum depth:} Introduced · \textbf{First seen:} Phase6 · \textbf{Latest development:} Phase6

\subsection{Definition and mechanism}\label{appendix-215ac7099a82-definition-and-mechanism}

Minkowski addition combines pairwise position coordinates of sets in a common frame. Collision difference means \(A+(-B)\) here. A support function returns an extreme scalar measurement; a support mapping returns a maximizing point, possibly with ties.

Intervals \([2,4]\) and \([3,5]\) have difference \([-3,1]\), containing zero. Support of the difference uses a support point of \(A\) minus one of \(B\) queried in the opposite direction.

\subsection{Prerequisites}\label{appendix-215ac7099a82-prerequisites}

\begin{itemize}
\tightlist
\item
  \hyperref[appendix-0026ec42c9d5]{Convex combinations and convex hulls} --- Convex set geometry is the domain for hull support and collision constructions.
\item
  \href{../06-GPU-Architecture-and-Real-Time-Rendering/main.pdf\#nameddest=appendix-8bc4e8fc547c}{Dot products and orthogonality (GPU Architecture and Real-Time Rendering, Reference notes: Dot products and orthogonality)} --- Inner-product measurements supply the perpendicularity and alignment constraints.
\end{itemize}

\subsection{Used by}\label{appendix-215ac7099a82-used-by}

Later curriculum sources can deepen this same node; no future coverage is marked complete.

\subsection{Recall and next depth}\label{appendix-215ac7099a82-recall-and-next-depth}

Prove the overlap/origin equivalence; separate the support identity from a complete GJK algorithm. Full production proofs, target-specific integration and learner mastery require further evidence.

\subsection{Source and navigation}\label{appendix-215ac7099a82-source-and-navigation}

\href{../02-Mathematics-for-Visual-Computing/main.pdf\#nameddest=01-foundations-i-ch018-weighted-points-lead-from-affine-geometry-to-collision-shapes}{Teaching narrative (Mathematics for Visual Computing, Mathematics for Visual Computing, Chapter 2)} · \hyperref[appendix-fd090d982a15-phase-6---parts-501-600]{Coverage and source gap} · \hyperref[appendix-a103e7846d21-phase-6---source-analysis-and-verification]{Evidence and limits} · \href{../01-Bits-Programs-and-Reliable-Systems/main.pdf\#nameddest=appendix-4ad47bdd269a}{Dependency Map (Bits, Programs, and Reliable Systems, Reference notes: Dependency Map)}
